\chapter{Inverse systems, Weyl--Heisenberg representation and predictive quotient}
\label{chap:numero-heisenberg-d108}
\index{HMT number}
\index{Heisenberg!finite group}

The construction distinguishes three levels: the finite seed catalogue, the
state space that preserves carry and orientation, and the inverse
limit of compatible prolongations. This stratification makes it possible to
formalize contraction of cylindrical intervals to a real point,
persistence of trajectory data in a single evaluation fiber, and
the finite Weyl--Heisenberg representation associated with the alternating form
of the nonadic torus supporting the additive and multiplicative observables.

The symbols \(V_t\), \(X_t\), and \(\rho_{t+1,t}\) used here are the
same sets of local states, admissible prefixes, and truncations
introduced in \cref{chap:numero-medicion}, with the purely
notational change \(t=m\) and \(\rho_{n,m}=p_{n,m}\). This section makes their
coordinates explicit and introduces no second inverse system.

\section{Observable projection and complete state space of the TPK}

The observable projection of the phase transport system (TPK) is the
finite encoding application
\[
104\,976\longrightarrow468\longrightarrow243.
\]
which classifies seeds, signatures, and words \(w_6\). The complete TPK
state space further preserves
\[
\begin{gathered}
\text{bidirectional orientation},\quad
\text{Hensel state},\quad
\text{cylinder and carry},\\
\text{joint boundary signature},\quad
\text{visible nonadic phase }g_{\rm vis}(m).
\end{gathered}
\]
Two candidates with the same local signature may belong to distinct classes
when the following boundary is conserved. We do not introduce a second tuple
of state. For
\[
x_m=(q_m,o_m,h_m,c_m,\sigma_m,C_m,b_m,\lambda_m)\in\XTPK^{(m)}
\]
we use the coordinate chart
\[
\chi_m(x_m)=
\bigl(
B_m,h_m,(C_m,b_m),\sigma_m,
o_m,\mu(q_m),g_{\rm vis}(m),\lambda_m
\bigr),
\]
with
\[
B_m=(u_m^\pi,u_m^e,u_m^\varphi)\in(\mathbb F_3^6)^3.
\]
Here \(h_m\) is the Hensel extension; \((C_m,b_m)\), the cylinder and its
boundary; \(o_m,\mu(q_m)\), the orientation and the bidirectional observable; and
\[
g_{\rm vis}(m)=1+((m-1)\bmod9)\in\{1,\ldots,9\},
\]
the visible nonadic phase label. The corresponding class coordinate
is \(g_0(\theta_m)\in\ZZ/9\ZZ\), with \(\theta_m=[m-1]_{108}\). The joint signature
is the coordinate already typed by the phase kernel,
\[
\sigma_m=\operatorname{Sig}_{q_m}(h_m,c_m,\lambda_m)
=(q_\partial,a_\partial,c_\partial,
\operatorname{colw},\rho_W,r_\partial).
\]
The joint signature does not reduce to three independent channel signatures, because the
selection associated with the involution \(\pi/e\) and the axis \(\varphi\) depends on
correlated data. Two nodes with the same \(B_m\) represent distinct
states when they differ in extension, cylinder, carry, signature, orientation,
or transport record.

\subsection{Exact compatibility of the cylinders under truncation and nonadic prolongation}

Let \(N\) be the integer of a ternary prefix of length \(L\), and let \(D\)
be the integer obtained by concatenating its \(b\) decimal triads. We define
\[
A=N1000^b-D3^L,
\qquad
B=(N+1)1000^b-D3^L.
\]
The cylinders intersect if and only if
\[
\boxed{A<3^L,\qquad B>0.}
\]
Adding a ternary block \(u\in\{0,\ldots,728\}\) gives
\[
N'=729N+u.
\]
If a new decimal triad does not appear,
\[
A'=729A+u1000^b.
\]
If \(\delta\) occurs, so that \(D'=1000D+\delta\), then
\[
A'=729000A+u1000^{b+1}-\delta\,729\,3^L.
\]
If \(\Delta b\in\{0,1\}\) records the appearance of that new triad, in
both cases
\[
B'=A'+1000^{b+\Delta b}.
\]
The transition, once the admissible pair \((u,\delta)\) has been fixed, does not consult a
target value. The joint filtration by intersection of cylinders and nonadic
return acts at a distinct level: it selects the boundary pairs that
admit compatible prolongation.

\subsection{Cyclic registration of 108 stages and its reduced projection}

We call the sequence of \emph{108-stage cyclic record} the sequence of
\(108\) steps grouped into twelve sectors of nine steps, with the index
convention fixed in \cref{chap:tpk-definicion}.

Let \(\XTPK^{[108]}\subseteq\XTPK\) be the domain of complete states whose
register \(\lambda\) contains the declared \(108\) steps. Let us denote by
\(\mathcal C_{108}^{\rm red}\) the set of registers that result by
conserving only the reduced fields, and by \(\mathcal U_{12}^{\rm sgn}\) the
set of admissible signed dodecaphases. The instance distinguishes two declared maps:
\[
F:\XTPK^{[108]}\longrightarrow\mathcal C_{108}^{\rm red},
\qquad
R:\XTPK^{[108]}\longrightarrow\mathcal U_{12}^{\rm sgn}.
\]
The map \(F\) retains only the fields of the reduced register, whereas \(R\) evaluates the signed dodecaphase. The existence of a map \(\Xi\) such that \(R=\Xi\circ F\) is equivalent to the factorization of \(R\) through a projection that may eliminate the observable coordinate, orientation, carry, or boundary. This factorization constitutes an additional proposition and is not an implicit part of the reconstruction on the complete state space.

The reduced record is, by definition, a projection of the complete state,
not the initial state. Subsequent results use \(R\) on the
complete state and do not presuppose a factorization through \(F\).

\section{Inverse systems of survivors}

The TPK transition need not be a function. To type this fact correctly,
let \(V_t\) be the finite set of possible complete states at
depth \(t\), and let
\[
\mathcal R_t\subseteq V_t\times V_{t+1}
\]
the admissible transition relation, with all fields of the
enriched state conserved. We then define \(X_t\), not as the set of
isolated terminal states, but as the set of admissible prefixes, that is,
of complete admissible histories:
\[
X_t=
\Bigl\{
(v_0,\ldots,v_t)\in\prod_{j=0}^{t}V_j:
(v_j,v_{j+1})\in\mathcal R_j\ (0\le j<t),
\ \text{and the closures are satisfied up to }t
\Bigr\}.
\]
On these objects, truncation is indeed a canonical application:
\[
\begin{aligned}
\rho_{t+1,t}:X_{t+1}&\longrightarrow X_t,\\
(v_0,\ldots,v_t,v_{t+1})&\longmapsto(v_0,\ldots,v_t).
\end{aligned}
\]
It is well defined precisely because every prefix of an admissible history remains an admissible history for the closures already imposed. If a filter were to use future information and could retroactively invalidate a prefix, \(X_t\) would first have to be replaced by the subset of prolongable prefixes; under that convention, the same truncation is again well defined.

A global HMT history is a compatible family of prefixes.
\[
x=(x_t)_{t\ge0},
\qquad
x_t\in X_t,
\qquad
\rho_{t+1,t}(x_{t+1})=x_t.
\]
It is equivalent to an infinite sequence \((v_0,v_1,\ldots)\) whose consecutive
pairs belong to the relations \(\mathcal R_t\) and whose prefixes
satisfy all closures. The history space is the inverse limit
\[
X_\infty=\varprojlim X_t.
\]
This formulation avoids imposing a single-valued local transition \(B_t\mapsto B_{t+1}\). The transition may be a finite relation, several branches may share the same visible terminal state, and uniqueness may appear only after future compatibility is imposed. By preserving the complete prefix, the information that makes it possible to distinguish them is not lost prematurely.

\begin{theorem}[Unique coinductive section]
Suppose that, for every \(t\), the survivor set \(S_t\subseteq
X_t\) is nonempty, the restrictions
\(\rho_{t+1,t}:S_{t+1}\to S_t\) are compatible, and every \(S_t\) contains a
unique prefix. Then \(\varprojlim S_t\) contains a unique global
history.
\end{theorem}

\begin{proof}
Let \(x_t\) be the unique prefix of \(S_t\). Compatibility forces
\(\rho_{t+1,t}(x_{t+1})=x_t\), so that \((x_t)_t\) belongs to the inverse
limit. Every other history would have a member of \(S_t\) at each level
and, by uniqueness, would coincide prefix by prefix with \((x_t)_t\).
\end{proof}

The theorem formalizes the passage from level-by-level compatibility to a global
history. It does not turn a finite window into a coinductive conclusion: its
application to a concrete family requires verifying at every depth non-
emptiness, compatibility of the restrictions, and uniqueness of the
surviving prefixes.

\section{Nonadic Survival Connection}
\label{sec:conexion-nonadica}

Cylinder contraction does not proceed block by block as a deterministic substitution. At certain boundaries, there are several local allocations of the same aggregate datum; prolongability determines which belongs to the published history. The resulting structure is a connection on the nine phase classes: upon completing one turn, the phase class returns, but the state of the fiber has changed. This section constructs the complete finite window in which that return is observed.

The domain of this local calculation is the transport register generated by the
APP--TRIT--TPK composition. Its initial states, its integer propagation
matrix and its boundary relations are prior internal publications of
the enriched transition \(\mathcal U_t\), in accordance with the oriented register of
twelve transitions and the unique reconstruction
\(L_i=X_i^{-1}Y_i\). This section receives them as already proved results,
not as external premises or conventional data. All its statements
are exact on those outputs. The preceding coinductive theorem extends the
same construction to arbitrary depth and preserves nonemptiness,
truncation compatibility and uniqueness of the surviving prefixes.

\subsection{Prefix of \(30\) trits and 1st boundary}

We write \(u_k^\chi\in\FF_3^6\), with
\(\chi\in\{\pi,e,\varphi\}\), for the visible block of channel \(\chi\) at
depth \(k\). The first five blocks form three prefixes of thirty
trits:
\begin{align*}
w_{30}^{\pi}&=010211|012222|010211|002111|110221,\\
w_{30}^{e}&=201101|121221|102011|012222|102011,\\
w_{30}^{\varphi}&=121200|112202|121020|010210|010200.
\end{align*}
The sixth block is the first boundary following those prefixes:
\[
R_{36}=
\begin{pmatrix}
u_5^\pi\\ u_5^e\\ u_5^\varphi
\end{pmatrix}
=
\begin{pmatrix}
222220\\021222\\102011
\end{pmatrix}.
\]
The notation \(R_{36}\) indicates that thirty-six trits per
channel have been displayed. It does not denote a terminal state: the record reconstructed below
contains twenty blocks, and the recurrence admits further continuations
when the corresponding Hensel states are supplied.

\subsection{Self-scale axis and residual fiber}

\ifdefined\HMTAxialTheoremDefined
The axial decomposition of the signature and the identification of the axis
\(\varphi\) against the residual fiber \(\pi/e\) were already proved in the
\cref{thm:eje-espejo-nonadico}. Here its consequence for the
nonadic window is conserved without duplicating the theorem or the proof.
\else
For
\[
B_k=(u_k^\pi,u_k^e,u_k^\varphi)\in(\FF_3^6)^3
\]
we define, coordinate by coordinate,
\[
q_k=u_k^\pi+u_k^e+u_k^\varphi,
\qquad
a_k=u_k^\pi+u_k^e-u_k^\varphi.
\]

\begin{theorem}[Axial decomposition of the boundary signature]
\label{thm:eje-espejo-nonadico}
The pair \((q_k,a_k)\) uniquely determines the axial channel and the aggregate of
the other two channels:
\[
u_k^\varphi=2(q_k-a_k),
\qquad
u_k^\pi+u_k^e=q_k-u_k^\varphi.
\]
Consequently, a fiber of the application
\[
q_{\rm ax}(u^\pi,u^e,u^\varphi)
=(u^\pi+u^e,u^\varphi)
\]
can preserve ambiguity only in the internal allocation between \(u^\pi\) and
\(u^e\).
\end{theorem}

\begin{proof}
Subtracting the two defining equations gives
\(q_k-a_k=2u_k^\varphi\). The inverse of \(2\) in \(\FF_3\) is again
\(2\), whence
\(u_k^\varphi=2(q_k-a_k)\). Subtracting this vector from \(q_k\) produces
\(u_k^\pi+u_k^e\). The operation does not recover each summand separately; for
this reason the unique residual fiber is the set of distributions of the fixed aggregate.
\end{proof}

In the model's geometric terminology, \(\varphi\) constitutes the axis
fixed by the signature and the pair \(\pi/e\) the residual fiber. The expression does not
postulate an analytic symmetry of the real constants: it names the
exact linear decomposition just proved in \(\FF_3^6\).
\fi

\subsection{Orbita complete of the \(9\) fases consecutivas of the Topological Phase Kernel}

We fix the absolute phase by
\[
g(k)=
\begin{cases}
k\bmod9,&k\not\equiv0\pmod9,\\
9,&k\equiv0\pmod9.
\end{cases}
\]
Since \(w_{30}\) occupies \(k=0,\ldots,4\), the first subsequent complete circuit
after the prefix traverses \(k=5,\ldots,13\), with phases
\(5,6,7,8,9,1,2,3,4\). Enumeration of the \(729=3^6\) candidates of
aggregate signature at each level produces the following window:
\begin{center}
\small
\begin{tabular}{@{}ccrrlll@{}}
\toprule
\(k\)&\(g(k)\)&locals&selected&\(u_k^\pi\)&\(u_k^e\)&\(u_k^\varphi\)\\
\midrule
5&5&1&1&222220&021222&102011\\
6&6&1&1&111201&201202&221021\\
7&7&3&1&212121&222210&200221\\
8&8&3&1&200121&212212&110110\\
9&9&2&1&100100&020112&010122\\
10&1&1&1&101222&112221&112002\\
11&2&1&1&022212&110001&100021\\
12&3&1&1&012012&202222&000120\\
13&4&1&1&111210&112102&211122\\
\bottomrule
\end{tabular}
\end{center}
The “local” column counts the admissible input states at the indicated
level; it does not count the outputs of the boundary relation. This distinction
is essential in \(k=9\): the two local input states precede a
three-output fiber of the selected state.

\subsection{Return 9→1 of the nonadic phase with state change and memory advance}

Let
\[
 \pi_B:V_k\longrightarrow\mathcal B:=(\FF_3^6)^3
\]
the projection onto the visible triple. Once the fiber data of the record are fixed,
the transport relation on complete states induces a visible
relation \(\MonNine^B\). The selected state entering the ninth phase is
\[
B_9=(100100\mid020112\mid010122).
\]
Its relational image consists exactly of
\[
\begin{aligned}
B_{10}^{(0)}&=(101222\mid112221\mid112002),\\
B_{10}^{(1)}&=(102221\mid111222\mid112002),\\
B_{10}^{(2)}&=(111221\mid102222\mid112002).
\end{aligned}
\]
The three outputs have
\[
q_{10}=022112,
\qquad a_{10}=101111,
\]
and, by \cref{thm:eje-espejo-nonadico}, share
\[
u_{10}^\varphi=112002,
\qquad u_{10}^\pi+u_{10}^e=210110.
\]
They differ exclusively in the three published distributions of that aggregate
between \(\pi\) and \(e\).

\begin{theorem}[Two-level extendibility selection]
\label{thm:seleccion-novena-fase}
In the calibrated window, the three preceding outputs admit a prolongation
of depth one. Upon also imposing the next block
\[
B_{11}=(022212\mid110001\mid100021)
\]
and the decimal boundary \((502,662,638)\), only \(B_{10}^{(0)}\) admits a
prolongation of the same complete state. Therefore
\[
\#\operatorname{Surv}^{\rm cal}_{1}(B_9)=3,
\qquad
\#\operatorname{Surv}^{\rm cal}_{2}(B_9)=1.
\]
\end{theorem}

\begin{proof}
The decimal triples associated with the three outputs are, respectively,
\[
(279,352,365),\quad(280,351,365),\quad(282,350,365).
\]
For a ternary word of length \(L\), with associated integer \(N\), and a
prefix of \(b\) triples, with associated integer \(D\), compatibility of the
two cylinders is equivalent to the two integer inequalities
\[
N1000^b<(D+1)3^L,
\qquad
(N+1)1000^b>D3^L.
\]
The six inequalities for each output are satisfied at the current level. Upon
adjoining \(B_{11}\) and \((502,662,638)\), they are again satisfied for the three
channels of \(B_{10}^{(0)}\). For \(B_{10}^{(1)}\) and \(B_{10}^{(2)}\), the
inequality of channel \(\varphi\) is preserved, but those of \(\pi\) and
\(e\) fail. The evaluation is performed with integers in the corresponding finite certificate; no
numerical tolerance intervenes. Consequently, the candidate set
is reduced from three to one when passing from depth one to depth two.
\end{proof}

The other two words can precede the same ternary tail only if the
deep Hensel extension is changed. They are therefore not prolongations of
the same complete state of the register: the visible residue partially coincides, but
the genealogy of the fiber is distinct.

\subsection{Phase return with state change}

The phase satisfies \(g(k+9)=g(k)\). The visible state, by contrast, is not
reset. In particular,
\[
B_1=(012222\mid121221\mid112202),
\qquad
B_{10}=(101222\mid112221\mid112002),
\]
so that \(g(1)=g(10)=1\) but \(B_1\ne B_{10}\). The verification of the
nine pairs \(k\mapsto k+9\), for \(k=1,\ldots,9\), also gives
\[
u_{k+9}^\varphi\ne u_k^\varphi,
\qquad
u_{k+9}^\pi+u_{k+9}^e\ne u_k^\pi+u_k^e.
\]
Therefore, the window satisfies exactly
\[
\boxed{g(k+9)=g(k),\qquad B_{k+9}\ne B_k
\quad(1\le k\le9).}
\]
This is a phase return with nontrivial visible holonomy. The designation
\emph{nonadic monodromy} refers to this return action on the
enriched fiber; it does not introduce a tenth phase. To promote the finite remark
to a global monodromy, the action at every depth must be constructed
and the coinductive compatibility required at the beginning of the section must be proved.

\subsection{The branch of \(20\) blocks}

The preceding return does not close the emission. The twenty blocks that the
certificate reconstructs from the three published integer states are
\begin{align*}
\pi:\;&010211|012222|010211|002111|110221|222220|111201|\\
&212121|200121|100100|101222|022212|012012|111210|\\
&121011|200220|120210|000101|022010|020111,\\[2pt]
e:\;&201101|121221|102011|012222|102011|021222|201202|\\
&222210|212212|020112|112221|110001|202222|112102|\\
&102010|022020|222002|012010|001112|022212,\\[2pt]
\varphi:\;&121200|112202|121020|010210|010200|102011|221021|\\
&200221|110110|010122|112002|100021|000120|211122|\\
&202200|202110|221000|100102|111122|020010.
\end{align*}
Thus, \(w_{30}\) is the prefix and \(R_{36}\) the first boundary of an
explicitly continued branch; neither object replaces the enriched state
that transports that continuation.

\subsection{Exact decoding of \(12\) decimal triples by means of the nonadic reader}

Let \(N_{\chi,k}\) be the integer represented in base three by the first \(k\)
blocks of channel \(\chi\). With \(k=13\), the word has \(78\) trits and
defines
\[
D_{\chi,12}
=\left\lfloor\frac{1000^{12}N_{\chi,13}}{3^{78}}\right\rfloor.
\]
The escritura of \(D_{\chi,12}\) in twelve blocks of three cifras is
\[
\begin{array}{c|rrrrrrrrrrrr}
\pi&141&592&653&589&793&238&462&643&383&279&502&884\\
e&718&281&828&459&045&235&360&287&471&352&662&497\\
\varphi&618&033&988&749&894&848&204&586&834&365&638&117.
\end{array}
\]
The identity is entire: it is equivalent to
\[
D_{\chi,12}3^{78}
\le 1000^{12}N_{\chi,13}
<(D_{\chi,12}+1)3^{78}.
\]
Therefore, the two cylinders intersect in each channel and the reading
operation produces exactly the thirty-six published digits. This
result is a consequence of the generated register. The historical
constructions that used conventional expansions belong exclusively
to comparison and subsequent checking: they do not enter into orbital
selection, the transition \(\mathcal U_t\), the reconstruction of
\(L_0,L_1\) or the calendar \(0011\). The identity certifies the Archimedean
reading of the states produced by APP--TRIT--TPK and preserves intact
their generative provenance independent of target values.

\subsection{Schedule for cylinder reopening and periodicity of the compatible prolongation}

The conversion of a six-trit block to decimal triples has average slope
\[
\lambda=6\log_{1000}3\in(0,1),
\qquad k(t)=\lfloor\lambda(t+1)\rfloor.
\]
An index \(t\) is a reopening without a new triple when
\(k(t+1)=k(t)\). Let \(\delta=1-\lambda\).

\begin{theorem}[Beatty Reading Calendar Ternary-Decimal]
\label{thm:calendario-beatty}
The set of reopenings is
\[
\mathcal T=
\left\{\left\lceil\frac{m}{\delta}\right\rceil-2:m\ge1\right\}.
\]
Its first elements are
\[
20,42,64,86,108,130,151,173,195,217,\ldots,
\]
and two reaperturas consecutivas distan \(21\) or \(22\) blocks.
\end{theorem}

\begin{proof}
The quantity \(\lambda=\log_{1000}729\) is irrational: an equality
\(\lambda=p/q\) would imply \(729^q=1000^p\), impossible by prime
factorization. Also \(\delta\) is irrational, so that
\[
k(t)=(t+1)-\lceil\delta(t+1)\rceil.
\]
The equality \(k(t+1)=k(t)\) occurs exactly when the ceiling function of the
second term increases by one. The \(m\)-th occurrence is obtained
by solving
\(\delta(t+1)<m<\delta(t+2)\), which gives
\(t=\lceil m/\delta\rceil-2\). The differences of a Beatty sequence
belong to
\[
\left\{\left\lfloor\delta^{-1}\right\rfloor,
\left\lceil\delta^{-1}\right\rceil\right\}=\{21,22\}.\qedhere
\]
\end{proof}

\subsection{Bockstein--Hensel Recurrence of Residue and Quotient}

The memory carried by the branch proceeds from the exact sequence.
\[
0\longrightarrow\ZZ^6\xrightarrow{\,\times3\,}\ZZ^6
\longrightarrow\FF_3^6\longrightarrow0
\]
and of the unimodular matrix
\[
A=
\begin{pmatrix}
2&2&2&1&2&1\\
2&1&2&2&1&1\\
1&0&0&1&0&2\\
0&0&1&1&1&0\\
2&2&2&0&2&1\\
2&1&2&1&0&0
\end{pmatrix},
\qquad\det A=1.
\]
For a row state \(x_k^\chi\in\ZZ^6\), it is defined
\[
y_k^\chi=x_k^\chi A,
\qquad
u_k^\chi=y_k^\chi\bmod3\in\{0,1,2\}^6,
\qquad
x_{k+1}^\chi=\frac{y_k^\chi-u_k^\chi}{3}.
\]

\begin{proposition}[Exact conservation of quotient memory]
\label{prop:memoria-bockstein-hensel}
With \(x_k^\chi\) and \(A\) fixed, the pair
\((u_k^\chi,x_{k+1}^\chi)\) is unique and satisfies
\[
x_k^\chi A=u_k^\chi+3x_{k+1}^\chi.
\]
Conversely, since \(A\in\operatorname{GL}_6(\ZZ)\), the pair uniquely determines
\(x_k^\chi\).
\end{proposition}

\begin{proof}
Euclidean division of each coordinate of \(x_k^\chi A\) by three produces
a unique residue in \(\{0,1,2\}\) and a unique integer quotient. This proves
the first assertion and the identity. Since \(\det A=1\), the inverse
\(A^{-1}\) has integer entries; therefore
\[
x_k^\chi=(u_k^\chi+3x_{k+1}^\chi)A^{-1}
\]
recovers the previous state unambiguously.
\end{proof}

Residue \(u_k^\chi\) is the visible part, and quotient
\(x_{k+1}^\chi\) the transported 3-adic memory. The certificate reproduces
sixty coordinated steps ---twenty per channel---, verifies \(AA^{-1}=I\) with
integer arithmetic, and checks all blocks against the register. The designation
Bockstein--Hensel refers here to this residue--quotient separation and its
3-adic iteration; it is not identified without additional data with a
cohomological Bockstein operator.

\section{From the prefix to a real point}

Let us associate with each compatible prefix \(x_t\in X_t\) a closed interval
\(I_t(x_t)\subset\mathbb R\). The geometric interpretation is that of a
cylinder: all futures sharing the same prefix project into
\(I_t\).

\begin{theorem}[Archimedean convergence of nested intervals]
Let \((x_t)_{t\ge0}\) be a compatible history. If
\[
I_{t+1}(x_{t+1})\subseteq I_t(x_t)
\]
and
\[
\operatorname{diam} I_t(x_t)\longrightarrow0,
\]
then there exists a unique \(r\in\mathbb R\) such that
\[
\bigcap_{t\ge0}I_t(x_t)=\{r\}.
\]
\end{theorem}

\begin{proof}
This is the principle of nested intervals. The left endpoints form a
bounded increasing sequence and the right endpoints a bounded decreasing one. Their limits
exist; the difference between them is bounded by the diameter, which tends to
zero. Both limits coincide at a single point contained in all the
intervals.
\end{proof}

Let \(X_\infty^{\rm Arch}\subseteq X_\infty\) be the domain of histories for
which the cylinders are nested and their diameters tend to zero. If those
two properties are part of global admissibility, then
\(X_\infty^{\rm Arch}=X_\infty\); if not, this domain restriction is
necessary. We define the evaluation
\[
\operatorname{ev}_{\mathbb R}:X_\infty^{\rm Arch}\longrightarrow\mathbb R,
\qquad
\operatorname{ev}_{\mathbb R}(x)=
\text{the unique point of }\bigcap_{t\ge0}I_t(x_t).
\]
Uniqueness of the value does not entail uniqueness of the history. Two
genealogies \(x\ne y\) may satisfy
\[
\operatorname{ev}_{\mathbb R}(x)
=\operatorname{ev}_{\mathbb R}(y).
\]
We define reading equivalence
\[
x\sim_{\rm ev}y
\quad\Longleftrightarrow\quad
\operatorname{ev}_{\mathbb R}(x)
=\operatorname{ev}_{\mathbb R}(y).
\]

\begin{theorem}[Archimedean evaluation quotient]
The map
\[
\begin{aligned}
\overline{\operatorname{ev}}_{\mathbb R}:
X_\infty^{\rm Arch}/\!\sim_{\rm ev}
&\longrightarrow
\operatorname{Im}(\operatorname{ev}_{\mathbb R})\subseteq\mathbb R,\\
[x]&\longmapsto\operatorname{ev}_{\mathbb R}(x)
\end{aligned}
\]
is a bijection. In particular, the quotient is identified with all of
\(\mathbb R\) if and only if the evaluation
\(\operatorname{ev}_{\mathbb R}\) is surjective.
\end{theorem}

\begin{proof}
The definition of \(\sim_{\rm ev}\) makes the value independent of the
representative, so the map is well defined. It is injective because
two classes with the same value are, by definition, a single class; it is
surjective onto the image by the very definition of image. The last
assertion is equivalent to
\(\operatorname{Im}(\operatorname{ev}_{\mathbb R})=\mathbb R\).
\end{proof}

The covering theorem for the Archimedean reader constructs precisely those
prefixes: for every compact interval \(K\subset\mathbb R\),
\(\operatorname{ev}_{\mathbb R}(\Omega_{\infty,K})=K\), and the intervals
\([-m,m]\) form a directed family whose union is \(\mathbb R\).
Consequently, every real number possesses at least one compatible HMT history, and the
proved assertion is
\[
X_\infty^{\rm Arch}/\!\sim_{\rm ev}
\cong\operatorname{Im}(\operatorname{ev}_{\mathbb R})
=\mathbb R.
\]
The Archimedean value is, therefore, the quotient that identifies the fibers of
the evaluation; the enriched HMT number also conserves the class of
trajectory, the carry, and the orientation. The identification does not follow from the
mere abstract formation of the quotient: it follows from the covering theorem already
proved in the discrete construction of the continuum.

Algebraic descent is another problem and must be solved operation by
operation. Suppose that two candidate operations
\(\oplus\) and \(\odot\) have been constructed on \(X_\infty^{\rm Arch}\) and are closed on that
domain. Addition descends to the quotient if and only if
\[
x\sim_{\rm ev}x',\ y\sim_{\rm ev}y'
\Longrightarrow
x\oplus y\sim_{\rm ev}x'\oplus y',
\]
while the product descends, independently, if and only if
\[
x\sim_{\rm ev}x',\ y\sim_{\rm ev}y'
\Longrightarrow
x\odot y\sim_{\rm ev}x'\odot y'.
\]
The first congruence does not imply the second. Moreover, for the induced operations
to be precisely the ordinary sum and product of the Archimedean quotient, it is necessary to prove separately
\[
\operatorname{ev}_{\mathbb R}(x\oplus y)
=\operatorname{ev}_{\mathbb R}(x)+\operatorname{ev}_{\mathbb R}(y),
\]
\[
\operatorname{ev}_{\mathbb R}(x\odot y)
=\operatorname{ev}_{\mathbb R}(x)\operatorname{ev}_{\mathbb R}(y),
\]
and that the image of the evaluation be closed under the corresponding operation.
Only after these certificates does the quotient inherit the structure of a
subring ---or of a subfield, if additive opposites and multiplicative inverses are also constructed---
of \(\mathbb R\). The set-theoretic bijection of the preceding theorem is formal;
the content specific to HMT lies in constructing the evaluation, characterizing its
fibers, proving coverage, and establishing each descent law.

\section{Weyl--Heisenberg Representation Associated with the Nonadic Torus}

The oriented area form of the nonadic torus is
\[
\sigma((a,b),(c,d))=ad-bc\pmod9.
\]
This alternating form is obtained, with the chosen orientation, by
integrating the mixed curvature of APP over parallelograms of the torus. Let
\(\omega=\exp(2\pi i/9)\) and
\(\mathcal H_9=\mathbb C^9\), with basis \(\{|n\rangle:n\in\mathbb Z_9\}\).
We define
\[
X|n\rangle=|n+1\rangle,
\qquad
Z|n\rangle=\omega^n|n\rangle.
\]
Then
\[
ZX=\omega XZ,
\]
and, more generally,
\[
(X^aZ^b)(X^cZ^d)
=\omega^{bc}X^{a+c}Z^{b+d}.
\]
Therefore
\[
 c\bigl((a,b),(c,d)\bigr)=bc\pmod9
\]
is the polarized \(2\)-cocycle that appears in this section of the central
extension. We explicitly define
\[
 \operatorname{Heis}_9
 =\{\omega^kX^aZ^b:a,b,k\in\ZZ/9\ZZ\}.
\]
It is a group of order \(9^3=729\). It is a central extension associated
with the translation group of the APP support \(X_9\), not a subgroup of the
graph or the APP path category\@.
The commutator of the two transports is
\[
(X^aZ^b)(X^cZ^d)(X^aZ^b)^{-1}(X^cZ^d)^{-1}
=\omega^{bc-ad}I.
\]
Consequently, its exponent is \(-\sigma((a,b),(c,d))\). The polarized cocycle
\(c\) and the form \(\sigma\) are not the same object: the second is,
under this convention, the opposite of the antisymmetrization
\(c(u,v)-c(v,u)\). The commutator phase is thus the inverse of the
Wilson phase for the fixed orientation.

To construct observables, we take
\[
Q_{\rm H}=\operatorname{diag}(0,1,\ldots,8),
\qquad
F_{jk}=\frac13\omega^{jk},
\qquad
P_{\rm H}=F^\dagger Q_{\rm H}F.
\]
The operators \(Q_{\rm H},P_{\rm H}\) are self-adjoint. Their eigenbases are mutually
unbiased because \(|F_{jk}|^2=1/9\). For every unit vector
\(\psi\in\mathcal H_9\), the Robertson inequality gives
\[
\operatorname{Var}_\psi(Q_{\rm H})\,
\operatorname{Var}_\psi(P_{\rm H})
\ge
\frac14\left|\langle\psi,[Q_{\rm H},P_{\rm H}]\psi\rangle\right|^2,
\]
and the entropy bound for mutually unbiased bases produces
\[
H_{Q_{\rm H}}(\psi)+H_{P_{\rm H}}(\psi)\ge\log9.
\]
Here \(H\) is the Shannon entropy, with natural logarithm, of the
corresponding spectral distributions. The constructed relations are
summarized, without attributing any additional physical causality to them, by
\[
\begin{gathered}
\text{alternating form }(X_9,\sigma)
\longrightarrow
\operatorname{Heis}_9,\\
\operatorname{Heis}_9\longrightarrow(X,Z)
\longrightarrow(Q_{\rm H},P_{\rm H})
\longrightarrow\text{finite uncertainty}.
\end{gathered}
\]
This is a finite algebraic model of operators; by itself it does not specify
physical dynamics or units of measurement.

\section{Cyclic dynamics on \texorpdfstring{\(\mathbb Z_9\times\mathbb Z_{12}\)}{Z9 x Z12} and minimal predictive quotient}

The finite dynamical system \(\mathfrak C_{9,12}\) makes it possible to determine the minimal
memory required for the observed evolution to be deterministic. Let
\[
X=\{(b,m):b\in\mathbb Z/9\mathbb Z, 0\le m<12\}
\]
with successor
\[
S(b,m)=
\begin{cases}
(b,m+1),&m<11,\\
(b+1\bmod9,0),&m=11.
\end{cases}
\]
The marginal labels are
\[
\beta(b,m)=4b\pmod{12},
\qquad
d(b,m)=1+(m+\beta(b,m)\pmod{12}).
\]
For a remark application \(q\), the inclusive future word of horizon \(h\) is
\[
Q_h^q(x)=\bigl(q(x),q(Sx),\ldots,q(S^hx)\bigr).
\]

\begin{theorem}[Minimum predictive quotient]
Let \(X\) finite, \(T:X\to X\) and \(q:X\to V\). Definamos
\[
\operatorname{Eq}(f):=\{(x,y)\in X^2:f(x)=f(y)\},
\qquad
E_h=\operatorname{Eq}(Q_h^q),
\qquad
E_\infty=\bigcap_{k\ge0}\operatorname{Eq}(q\circ T^k).
\]
Then
\(E_{h+1}=E_h\cap(T\times T)^{-1}(E_h)\), the sequence stabilizes in finite
time, and \(X/E_\infty\) is the quotient with the smallest number of states that
refines \(q\) and admits an induced deterministic dynamics.
\end{theorem}

\begin{proof}
\(E_h\) requires equality of observations at the times \(0,\ldots,h\), whereas
\((T\times T)^{-1}(E_h)\) requires equality at \(1,\ldots,h+1\). Their intersection is
\(E_{h+1}\). Because \(X\) is finite, the number of classes cannot grow
indefinitely. At the first stable level, \(T\) respects the classes and
descends. Every \(T\)-compatible equivalence contained in
\(\operatorname{Eq}(q)\)
identifies only states with the same complete future and is therefore
contained in \(E_\infty\); hence the minimality property.
\end{proof}

The evaluation exhaustiva of \(\mathfrak C_{9,12}\) da:
\begin{center}
\small
\begin{tabular}{@{}lrrrr@{}}
\toprule
observable & instantaneous values & index \(h\) & observations & stable quotient\\
\midrule
\(\beta\) & 3 & 11 & 12 & \(C_{36}\), fibra 3\\
\(d\) & 12 & 8 & 9 & \(C_{36}\), fibra 3\\
\((\beta,d)\) & 36 & 0 & 1 & \(C_{36}\), fibra 3\\
\bottomrule
\end{tabular}
\end{center}
In the two marginal observables,
\[
\operatorname{Eq}(Q_{11}^{\beta})
=\operatorname{Eq}(Q_8^d)
=\operatorname{Eq}(\pi_{36}),
\qquad
\pi_{36}(b,m)=(b\bmod3,m).
\]
The factor \(36\) is not introduced as a target. It appears as the minimal predictive quotient of each marginal. The phase \(m\) is a state coordinate required to predict those futures; this mathematical assertion does not presuppose a physical causal interpretation. The transformation \(b\mapsto b+3\) is a symmetry invisible to those observables at every time.

As an abstract dynamical system, the quotient is a cycle of thirty-six states. Writing \(C_{36}\) does not, however, choose a distinguished origin: it must be fixed to obtain a coordinate \(\mathbb Z/36\mathbb Z\). Nor does this automatically turn those states into the thirty-six pairs of a nonad. Indeed, \(\operatorname{Aut}J(9,2)\cong S_9\). If the two points of a pair lie in distinct cycles of lengths \(r,s\) of a permutation, the orbit of the pair has length \(\operatorname{mcm}(r,s)\), with \(r+s\leq9\), and hence length at most \(20\). If they lie in the same cycle of length \(r\leq9\), the orbit has length \(r\), except in the antipodal case, when it has length \(r/2\). Thus no automorphism induces an orbit of length \(36\) on the pairs. Therefore, the passage
\[
C_{36}\longleftrightarrow \binom{[9]}2
\]
remains a calibrated identification, not a natural equivariant identification.
Even a Hamiltonian ordering of the pairs would prove local preservation
of incidence step by step, not a global automorphism or the canonicity of the
dictionary.

\section{HMT number as a compatible history in \(X_\infty\): genealogical identity,
reading structure, and Archimedean evaluation}

The preceding results determine the following instance of the governing definition.
In the notation of this section, an HMT number is a compatible
history
\[
x\in X_\infty=\varprojlim X_t.
\]
A reader \(L\) is additional structure: the pair \((x,L)\) constitutes a read presentation or a measurement, but does not alter the identity type of \(x\). When \(x\) belongs to \(X_\infty^{\rm Arch}\), its Archimedean value is the point intersection of the cylinders; the set of values so constructed is \(\operatorname{Im}(\operatorname{ev}_{\mathbb R})\subseteq\mathbb R\), and its identification with the entire line requires the additional coverage theorem. Its minimal memory relative to an observable is the predictive quotient of futures; and the alternating form of the sum--product support admits the finite Weyl--Heisenberg representation constructed above. Value, trajectory class, and uncertainty are defined by three operations on specified domains of the same transport architecture.
